Kesimpulan dalam proyek perancangan teknik tidak hanya berfungsi sebagai ringkasan, tetapi sebagai penegasan bahwa permasalahan teknik yang dirumuskan pada awal penelitian telah diselesaikan secara terukur dan dapat dipertanggungjawabkan secara akademis maupun teknis. Oleh karena itu, bagian ini disusun secara sistematis mencakup evaluasi tujuan, ringkasan desain, hasil pengujian, keterbatasan, serta rencana pengembangan lanjutan.

\subsection{Evaluasi Ketercapaian Tujuan}
\label{subsec:evaluasi_tujuan}

Permasalahan utama yang diangkat dalam proyek ini adalah \textit{[tuliskan permasalahan utama secara singkat]}. Berdasarkan hasil perancangan dan pengujian yang telah dilakukan, dapat disimpulkan bahwa tujuan utama proyek, yaitu \textit{[tuliskan tujuan utama]}, telah \textbf{[tercapai / sebagian tercapai / belum tercapai]}.

Pencapaian ini ditunjukkan melalui \textit{[indikator utama keberhasilan, misalnya: performa sistem, keberhasilan implementasi fitur, atau stabilitas operasi]}. Dengan demikian, solusi yang dirancang dapat dikatakan \textit{[efektif / cukup efektif / perlu pengembangan lebih lanjut]} dalam menjawab permasalahan yang telah dirumuskan.

\subsection{Ringkasan Desain Akhir}
\label{subsec:ringkasan_desain}

Desain akhir yang dihasilkan dalam proyek ini berupa \textit{[deskripsikan sistem/produk secara singkat, misalnya: sistem kontrol, perangkat lunak, alat berbasis sensor, dll.]}. Sistem ini terdiri dari beberapa komponen utama, yaitu:
\begin{itemize}
    \item \textit{[Komponen 1, misalnya: sensor / modul input]}
    \item \textit{[Komponen 2, misalnya: unit pemrosesan / mikrokontroler]}
    \item \textit{[Komponen 3, misalnya: aktuator / output sistem]}
    \item \textit{[Komponen 4, jika ada]}
\end{itemize}

Secara arsitektural, sistem bekerja dengan prinsip \textit{[misalnya: kontrol tertutup, pemrosesan berbasis sinyal, atau pipeline data]}, di mana sinyal masukan diproses dan menghasilkan keluaran yang sesuai dengan spesifikasi desain.

\subsection{Hasil Pengujian dan Validasi Kuantitatif}
\label{subsec:hasil_pengujian}

Berdasarkan hasil pengujian yang telah dilakukan, sistem menunjukkan kinerja sebagai berikut:

\begin{itemize}
    \item Waktu respons sistem rata-rata sebesar \textbf{[nilai]} detik, memenuhi spesifikasi awal yaitu maksimal \textbf{[nilai]} detik.
    \item Tingkat akurasi sistem mencapai \textbf{[nilai]}$\%$, dibandingkan dengan target awal sebesar \textbf{[nilai]}$\%$.
    \item Sistem mampu beroperasi secara stabil selama \textbf{[durasi]} tanpa mengalami kegagalan signifikan.
\end{itemize}

Hasil tersebut menunjukkan bahwa sistem yang dirancang telah memenuhi sebagian besar kriteria performa yang ditetapkan pada tahap perancangan awal.

\subsection{Keterbatasan Desain}
\label{subsec:keterbatasan}

Meskipun sistem telah menunjukkan kinerja yang baik, terdapat beberapa keterbatasan yang perlu diperhatikan, antara lain:

\begin{itemize}
    \item Sistem hanya beroperasi optimal pada kondisi \textit{[misalnya: suhu, lingkungan, atau rentang input tertentu]}.
    \item Keterbatasan pada \textit{[misalnya: kapasitas komputasi, resolusi sensor, atau bandwidth]} yang mempengaruhi performa sistem.
    \item Faktor eksternal seperti \textit{[noise, gangguan sinyal, atau keterbatasan perangkat keras]} yang belum sepenuhnya teratasi.
\end{itemize}

Keterbatasan ini menjadi pertimbangan penting dalam pengembangan sistem lebih lanjut.

\subsection{Rencana Implementasi dan Rekomendasi Pengembangan}
\label{subsec:future_work}

Sebagai tindak lanjut dari proyek ini, terdapat beberapa rekomendasi yang dapat dilakukan untuk pengembangan lebih lanjut:

\begin{itemize}
    \item Mengoptimalkan algoritme \textit{[misalnya: kontrol, filtering, atau komputasi]} untuk meningkatkan efisiensi dan performa sistem.
    \item Menggunakan komponen dengan spesifikasi lebih tinggi, seperti \textit{[sensor atau prosesor]}, untuk meningkatkan akurasi dan keandalan.
    \item Melakukan pengujian lanjutan dalam kondisi lingkungan yang lebih beragam untuk meningkatkan robustnes sistem.
    \item Mengembangkan sistem ke tahap implementasi nyata atau integrasi dengan sistem lain yang lebih kompleks.
\end{itemize}

Dengan adanya pengembangan tersebut, diharapkan sistem dapat memiliki performa yang lebih optimal serta dapat diimplementasikan secara lebih luas dalam aplikasi nyata.